\documentclass{article}
\usepackage{amsmath}


\begin{document}




Well\ldots this is a test to be honest, yes I am using a highly convinient typsetting language to shitpost but listen me for a
second bruh when math fuckery graduates from Basic-C-Repo (Which will in the future) math must be in a readable form + I need to FUCKING practice LaTeX duh\ldots


So here is the formula (actual math this time):
\begin{align*}
    \sin x &= \sum_{n=0}^{\infty}{(-1)}^n\frac{x^{2n+1}}{(2n+1)!}
            = x - \frac{x^3}{3!} + \frac{x^5}{5!} - \cdots, \\
    \cos x &= \sum_{n=0}^{\infty}{(-1)}^n\frac{x^{2n}}{(2n)!}
            = 1 - \frac{x^2}{2!} + \frac{x^4}{4!} - \cdots.
\end{align*}

Technically those are Maclaurin series, which is just Taylor series centered
at zero. Inputs are in radians (duh\ldots). Also that exclamation mark is a
factorial, the equation isn't yelling at you yet.

Now how does this become \texttt{cos\_and\_sinThingy.c}? Because writing
infinity into a for loop sounds like a great way to never see my donut again.

\section*{First make the angle fucking chill}

The series converges for every real angle, but throwing a huge number directly
into those powers is asking for floating-point fuckery. Sine and cosine repeat
every $2\pi$, so \texttt{reduce\_angle} first brings the input into
$[-\pi,\pi]$ without changing which point on the circle we mean.

If the angle is already in that interval, it just returns it. Otherwise the
code takes its magnitude and starts with \texttt{period} equal to $2\pi$.
It doubles that period while
\[
    \texttt{period} \leq \frac{|x|}{2},
\]
then walks back down by halving it, subtracting each period if it fits inside
the remaining magnitude. Basically subtract big chunks first instead of
subtracting $2\pi$ a billion times like an idiot (me, probably). The division
by two in that condition also keeps the next doubling from overflowing.

After that, the remainder satisfies $0 \leq r < 2\pi$ in exact arithmetic. If it is bigger
than $\pi$, subtract another $2\pi$; finally restore the original input's sign.
For example, $3\pi/2$ becomes $-\pi/2$. Same point on the circle, smaller
number to fuck with.

Then both functions fold the angle again, this time into
$[-\pi/2,\pi/2]$. Call the reduced angle $r$ and the folded angle $u$:
\[
    u =
    \begin{cases}
        \pi-r,  & r > \pi/2, \\
        -\pi-r, & r < -\pi/2, \\
        r,      & \text{otherwise}.
    \end{cases}
\]
Sine keeps its value under those folds: $\sin r = \sin u$.
Cosine needs a sign flip when either fold happens, so
\texttt{taylor\_cos} starts with \texttt{sign = 1.0} and changes it to
$-1$ in those two branches. That's why its return is
\texttt{sign * sum}. Smaller angle, same answer, fewer numerical headaches.

\section*{No factorial factory needed}

I could calculate every power and factorial separately, but the next term
already knows how to come from the previous one. For sine, divide two
consecutive terms and everything cancels down to
\[
    t_0=u, \qquad
    t_n=-t_{n-1}\frac{u^2}{(2n)(2n+1)}.
\]
For cosine it is the same trick with the even powers:
\[
    t_0=1, \qquad
    t_n=-t_{n-1}\frac{u^2}{(2n-1)(2n)}.
\]
That's literally what the two loops do. \texttt{x\_squared} stores $u^2$
once, \texttt{term} holds the current term, and \texttt{sum} collects them.
The minus sign alternates the terms for us. No \texttt{pow}, no factorial
function, no \texttt{math.h}; just multiplication, division and addition.

Both loops run from $n=1$ through $n=10$, but the initial term is already
in \texttt{sum}, so we actually use eleven terms:
\begin{align*}
    \sin u &\approx \sum_{n=0}^{10}{(-1)}^n\frac{u^{2n+1}}{(2n+1)!}
             &&(\text{through }-u^{19}/19! + u^{21}/21!), \\
    \cos u &\approx \sum_{n=0}^{10}{(-1)}^n\frac{u^{2n}}{(2n)!}
             &&(\text{through }-u^{18}/18! + u^{20}/20!).
\end{align*}
The code uses a fixed number of terms, not a ``keep going until accurate
enough'' check. Since $|u|\leq\pi/2$, the omitted alternating tail has shrinking
terms, and its first term bounds the truncation error:
\[
    |E_{\sin}|\leq\frac{{(\pi/2)}^{23}}{23!}, \qquad
    |E_{\cos}|\leq\frac{{(\pi/2)}^{22}}{22!}.
\]
Those are bounds for the series truncation in exact arithmetic, not a promise
that my C doubles magically have infinite precision. Rounding still exists,
because of course it does.

\section*{The weird inputs and the catch}

\verb|x != x| catches NaN and returns it. Comparing against
\texttt{DBL\_MAX} catches positive and negative infinity, and
\verb|x - x| turns those into NaN under the usual IEEE floating-point
behavior. There isn't a useful angle hiding inside infinity, bruh.
The fallback \texttt{\_\_DBL\_MAX\_\_} is a compiler-provided macro;
the standard header for \texttt{DBL\_MAX} is \texttt{<float.h>}.

Sine also returns immediately when \texttt{x == 0.0}. That preserves negative
zero too. Yes, negative zero is a thing, apparently regular zero wasn't enough.
Cosine at zero starts at $1$ and all the remaining terms are zero, as expected.

The big catch is angle reduction: \texttt{PI} and \texttt{TWO\_PI} are rounded
floating-point constants. For very large angles, reducing with that rounded
period can give a bad remainder, and doubles themselves stop resolving tiny
changes between huge numbers. Adding more Taylor terms won't fix an angle
that was already wrong before the loop started.

So for the spinning ASCII donut, keep the rotation angles bounded instead of
letting them grow forever (yes, future me, this part is for you). This is my
small sin/cos implementation for the experiment, not a claim that I somehow
outsmarted the entire math library in two for loops lol.



\end{document}
